% WAFC -- suplemento, versao k = 1 (E5a). Template da Statistica Sinica
% (manuscript/ss-template/supp-temp_20240820.tex). Traz as provas dos
% resultados de ms_1.tex e os lemas auxiliares. As fontes sao
% derivations/01 a 05; o comentario acima de cada secao diz qual.
\documentclass[10pt]{book}
\usepackage[sectionbib]{natbib}
\usepackage{array,epsfig,fancyheadings,rotating}
\usepackage{hyperref}
%%%%%%%%%%%%%%%%%%%%%%%%%%%%%%%%%%%%%%%%%%%%%%%%%%%%%%%%%%%%%%%%%%%%%%%%%%%%%%

\textwidth=31.9pc
\textheight=46.5pc
\oddsidemargin=1pc
\evensidemargin=1pc
\headsep=15pt
\topmargin=.6cm
\parindent=1.7pc
\parskip=0pt

\usepackage{amsmath}
\usepackage{amssymb}
\usepackage{amsfonts}
\usepackage{amsthm}
\usepackage{enumitem}
\usepackage{xcolor}
\usepackage[normalem]{ulem}
\definecolor{colR1}{RGB}{0,90,200}

\setcounter{page}{1}
\newtheorem{theorem}{Theorem}[section]
\newtheorem{lemma}{Lemma}[section]
\newtheorem{corollary}{Corollary}[section]
\newtheorem{proposition}{Proposition}[section]
\theoremstyle{definition}
\newtheorem{definition}{Definition}[section]
\newtheorem{example}{Example}[section]
\newtheorem{remark}{Remark}[section]
\pagestyle{fancy}

% --- macros: as mesmas de ms_1.tex (docs/notacao.md) ---------------------
\newcommand{\R}{\mathbb{R}}
\newcommand{\E}{\mathrm{E}}
\newcommand{\Prob}{\mathrm{P}}
\newcommand{\Var}{\mathrm{Var}}
\newcommand{\T}{^{\top}}
\newcommand{\norm}[1]{\left\lVert #1 \right\rVert}
\newcommand{\normn}[1]{\left\lVert #1 \right\rVert_{n}}
\newcommand{\inner}[2]{\langle #1, #2 \rangle}
\newcommand{\cv}{\ell}
\newcommand{\md}{m}
\newcommand{\lev}{j}
\newcommand{\tsl}{k}
\newcommand{\jzero}{j_{0}}
\newcommand{\NJ}{N_{J}}
\newcommand{\bX}{\mathbf{X}}
\newcommand{\bU}{\mathbf{U}}
\newcommand{\bZ}{\mathbf{Z}}
\newcommand{\fcoef}[1]{\beta_{#1}}
\newcommand{\lvl}[1]{c_{#1}}
\newcommand{\gcomp}[2]{g_{#1#2}}
\newcommand{\wav}[2]{\psi_{#1#2}}
\newcommand{\VJ}{V_{J}}
\newcommand{\WJ}{W_{J}}
\newcommand{\PiJ}{\Pi_{J}}
\newcommand{\coef}[4]{\theta_{#1#2,#3#4}}
\newcommand{\coefs}[4]{\theta^{*}_{#1#2,#3#4}}
\newcommand{\btheta}{\boldsymbol{\theta}}
\newcommand{\thetastar}{\btheta^{*}}
\newcommand{\thetahat}{\widehat{\btheta}}
\newcommand{\bbeta}{\boldsymbol{\beta}}
\newcommand{\betastar}{\bbeta^{*}}
\newcommand{\betahat}{\widehat{\bbeta}}
\newcommand{\bc}{\mathbf{c}}
\newcommand{\chat}{\widehat{\mathbf{c}}}
\newcommand{\Gram}{\Sigma}
\newcommand{\Gramhat}{\widehat{\Sigma}}
\newcommand{\pen}{\lambda}
\newcommand{\penn}{\lambda_{n}}
\newcommand{\Szero}{S_{0}}
\newcommand{\Besov}[3]{B^{#1}_{#2,#3}}
\newcommand{\seff}{s'}
\newcommand{\Amat}{\mathbf{A}}
\newcommand{\Bmat}{\mathbf{B}}
\newcommand{\Btil}{\widetilde{\Bmat}}
\newcommand{\PA}{\Pi_{\Amat}}
\newcommand{\MA}{M_{\Amat}}
\newcommand{\bveps}{\boldsymbol{\varepsilon}}
\newcommand{\bb}{\mathbf{b}}
\newcommand{\bY}{\mathbf{Y}}
\newcommand{\bPsi}{\boldsymbol{\Psi}}
\newcommand{\smax}{\widehat{\sigma}_{\max}}
\newcommand{\gtil}{\widetilde{\gamma}}
\newcommand{\Mstar}{M^{\star}}
\newcommand{\lmin}{\lambda_{\min}}
\newcommand{\lmaxx}{\lambda_{\max}}
\newcommand{\opnorm}[1]{\left\lVert #1 \right\rVert_{\mathrm{op}}}
\newcommand{\compat}{\phi_{0}}
\newcommand{\GBA}{\Gramhat_{\Bmat\mid\Amat}}
\newcommand{\thbar}{\bar{\btheta}}
\newcommand{\bbar}{\bar{\bb}}
\newcommand{\thstarM}[1]{\btheta^{*,(#1)}}
\DeclareMathOperator{\vecop}{vec}
\DeclareMathOperator{\tr}{tr}
\DeclareMathOperator{\supp}{supp}

%%%%%%%%%%%%%%%%%%%%%%%%%%%%%%%%%%%%%%%%%%%%%%%%%%%%%%%%%%%%%%%%%%%%%%%%%%%%%%
\pagestyle{fancy}
\def\n{\noindent}
\lhead[\fancyplain{} \leftmark]{}
\chead[]{}
\rhead[]{\fancyplain{}\rightmark}
\cfoot{}

%%%%%%%%%%%%%%%%%%%%%%%%%%%%%%%%%%%%%%%%%%%%%%%%%%%%%%%%%%%%%%%%%%%%%%%%%%%%%%
\begin{document}

\renewcommand{\baselinestretch}{2}

\markright{ \hbox{\footnotesize\rm Statistica Sinica: Supplement
}\hfill\\[-13pt]
\hbox{\footnotesize\rm
}\hfill }

\markboth{\hfill{\footnotesize\rm AUTHOR(S)} \hfill}
{\hfill {\footnotesize\rm WAVELET ADDITIVE FUNCTIONAL COEFFICIENTS} \hfill}

\renewcommand{\thefootnote}{}
$\ $\par \fontsize{12}{14pt plus.8pt minus .6pt}\selectfont

%%%%%%%%%%%%%%%%%%%%%%%%%%%%%%%%%%%%%%%%%%%%%%%%%%%%%%%%%%%%%%%%%%%%%%%%%%%%%%

\centerline{\large\bf SPARSE WAVELET ESTIMATION OF}
\vspace{2pt}
\centerline{\large\bf ADDITIVE FUNCTIONAL COEFFICIENTS}
\vspace{.25cm}
\centerline{Author(s)}
\vspace{.4cm}
\centerline{\it Affiliation(s)}
\vspace{.55cm}
\centerline{\bf Supplementary Material}
\vspace{.55cm}
\fontsize{9}{11.5pt plus.8pt minus .6pt}\selectfont
\noindent
This supplement contains the proofs of the results stated in the paper and
the auxiliary lemmas they use. Section S1 fixes the notation and the objects.
Section S2 proves the identification results of Section 2 of the paper,
Section S3 the approximation lemma, Section S4 the bounds for the Gram matrix
of the product design, Section S5 the oracle inequality, and Section S6 the
rates. References to numbered results without a section prefix are to the
paper; results numbered with the prefix S are stated and proved here.
\par

\setcounter{section}{0}
\setcounter{equation}{0}
\def\theequation{S\arabic{section}.\arabic{equation}}
\def\thesection{S\arabic{section}}

\fontsize{12}{14pt plus.8pt minus .6pt}\selectfont

%%%%%%%%%%%%%%%%%%%%%%%%%%%%%%%%%%%%%%%%%%%%%%%%%%%%%%%%%%%%%%%%%%%%%%%%%%%%%%
\section{Notation and objects}

The model, the basis and the estimator are those of Sections 2.1 to 2.3 of the
paper. We write $\bZ = [\Amat\;\Bmat]$ for the design matrix of equation
(2.4), with $\Amat \in \R^{n\times p}$ the unpenalized block and
$\Bmat \in \R^{n\times d}$, $d = pq\NJ$, the block of products; $\bY$, $f$ and
$\bveps$ for the vectors with entries $Y_i$, $f(\bX_i,\bU_i)$ and
$\varepsilon_i$; $\Gramhat = \bZ\T\bZ/n$ and $\Gram = \E(\bZ\bZ\T)$; and, for
$a, h \in \R^n$, $\normn{a}^2 = n^{-1}\sum_i a_i^2$ and
$\inner{a}{h}_n = n^{-1}\sum_i a_ih_i$.

The sieve oracle is $\betastar = (\bc,\thetastar)$, with
$\coefs{\cv}{\md}{\lev}{\tsl} = \inner{\gcomp{\cv}{\md}}{\wav{\lev}{\tsl}}$
for $\lev < J$, so that $f_J = \bZ\betastar$ is the best approximation of $f$
in the sieve in the sense of Lemma 2; the deterministic residual is
$\bb = f - f_J$, the oracle support is $\Szero = \supp(\thetastar)$ and
$s_0 = |\Szero|$, and $v = \thetahat - \thetastar$. When $\Amat$ has rank $p$
we write $\PA = \Amat(\Amat\T\Amat)^{-1}\Amat\T$, $\MA = I_n - \PA$,
$\Btil = \MA\Bmat$, $\GBA = \Btil\T\Btil/n$ and $\gtil = \lmin(\GBA)$, and
finally $\smax^2 = \max_a(\Gramhat)_{aa}$, the maximum being over the $d$
penalized columns.

For a support $S$ and a symmetric matrix $A$, the compatibility constant is
\[
  \compat^2(S;A) = \min\bigl\{|S|\,\delta\T A \delta :\;
  \norm{\delta_S}_1 = 1,\; \norm{\delta_{S^c}}_1 \le 3 \bigr\},
\]
and the elementary bound
$\delta\T A\delta \ge \lmin(A)\norm{\delta}_2^2 \ge
\lmin(A)\norm{\delta_S}_1^2/|S|$ gives
\begin{equation}
  \label{eq:compat-lmin}
  \compat^2(S;A) \;\ge\; \lmin(A) \qquad \text{for every } S,
\end{equation}
so that a lower bound on the smallest eigenvalue is a compatibility bound on
every support, with the same constant and no cone restriction.

%%%%%%%%%%%%%%%%%%%%%%%%%%%%%%%%%%%%%%%%%%%%%%%%%%%%%%%%%%%%%%%%%%%%%%%%%%%%%%
% derivations/01-identificabilidade.md
\section{Identification and the periodized basis}
\setcounter{equation}{0}

\subsection*{Proof of Proposition 1}

\emph{Levels and components are determined by the coefficients.} Let $\beta$
and $\tilde\beta$ be measurable maps from $[0,1]^q$ to $\R^p$ with
$\sum_{\cv}\fcoef{\cv}(\bU)X_{\cv} = \sum_{\cv}\tilde\beta_{\cv}(\bU)X_{\cv}$
almost surely, and put $\delta(u) = \beta(u)-\tilde\beta(u)$. Then
$\{\delta(\bU)\T\bX\}^2 = 0$ almost surely, and since $\delta(\bU)$ is
measurable with respect to $\bU$ and $\E\norm{\bX}_2^2 < \infty$,
\[
  0 \;=\; \E\bigl[\{\delta(\bU)\T\bX\}^2 \mid \bU\bigr]
  \;=\; \delta(\bU)\T\,\E(\bX\bX\T\mid\bU)\,\delta(\bU)
  \qquad \text{almost surely.}
\]
The conditional second moment matrix is positive definite almost surely, and a
positive definite quadratic form vanishes only at the origin, so
$\delta(\bU) = 0$ almost surely.

\emph{Levels and components are determined within a coefficient.} Because the
density of $\bU$ is positive Lebesgue almost everywhere, a set is null for the
law of $\bU$ if and only if it is Lebesgue null, so the identity
$\lvl{\cv} + \sum_{\md}\gcomp{\cv}{\md}(u_{\md}) =
\tilde{c}_{\cv} + \sum_{\md}\tilde{g}_{\cv\md}(u_{\md})$ holds Lebesgue almost
everywhere on $[0,1]^q$. Write $\Delta = \lvl{\cv}-\tilde{c}_{\cv}$ and
$\delta_{\md} = \gcomp{\cv}{\md}-\tilde{g}_{\cv\md} \in L_2[0,1] \subset
L_1[0,1]$, so that $\Delta + \sum_{\md}\delta_{\md}(u_{\md}) = 0$ almost
everywhere. Each summand is integrable on the cube by Tonelli's theorem.
Integrating in all coordinates and using $\int_0^1\delta_{\md} = 0$, which
holds because both parametrizations satisfy (2.2), gives $\Delta = 0$; fixing
$\md$ and integrating in the remaining coordinates gives
$\delta_{\md}(u_{\md}) = 0$ for almost every $u_{\md}$. \qed

\begin{remark}
\label{rem:centering}
If the centering $\int_0^1 \gcomp{\cv}{\md} = 0$ is replaced by
$\E\{\gcomp{\cv}{\md}(U_{\md})\} = 0$, the same proof applies under the
two sided density bound of Assumption 3, since then $L_2[0,1]$ and
$L_2(P_{U_{\md}})$ coincide as sets: the step that uses Fubini's theorem
yields $\delta_{\md}$ constant almost everywhere, and the constant is its
mean, which is zero. The two parametrizations are related by
$\lvl{\cv} \mapsto \lvl{\cv} + \sum_{\md}\E\{\gcomp{\cv}{\md}(U_{\md})\}$, as
stated in Section 2.1 of the paper. The difference is not cosmetic for the
estimator: the design of Section 2.3, which omits the constant scaling
function, estimates the version centered with respect to Lebesgue measure.
\end{remark}

\subsection*{Proof of Lemma 1}

\emph{The basis.} The periodization of the scaling function is
$\phi^{\rm per}_{00}(u) = \sum_{\iota}\phi(u+\iota)$, a finite sum at each $u$
because the support is compact, and it equals the constant function one, since
the integer translates of an orthonormal scaling function with
$\int\phi = 1$ form a partition of unity. For a wavelet, exchanging the finite
sum with the integral,
$\int_0^1\wav{\lev}{\tsl} = \int_{\R}2^{\lev/2}\psi(2^{\lev}y-\tsl)\,dy =
2^{-\lev/2}\int_{\R}\psi = 0$ by the vanishing moment of order zero. That
$\{\mathbf{1}\}\cup\{\wav{\lev}{\tsl}\}$ is an orthonormal basis of $L_2[0,1]$
is the standard construction of the periodized multiresolution analysis
\citep[Section 9.3]{Daubechies-1992}. Since $\dim\VJ = 2^J$ and the $2^J-1$
wavelets of level below $J$ are orthogonal to the constant, $\WJ$ is exactly
the set of elements of $\VJ$ with zero integral.

\emph{Injectivity.} Suppose two parameter vectors give the same regression
function almost surely. The functions
$\beta^J_{\cv}(u) = \lvl{\cv}+\sum_{\md}g^J_{\cv\md}(u_{\md})$, with
$g^J_{\cv\md} = \sum_{\lev\tsl}\coef{\cv}{\md}{\lev}{\tsl}\wav{\lev}{\tsl}$,
are measurable and bounded, so the first part of Proposition 1 gives
$\beta^J_{\cv} = \tilde\beta^J_{\cv}$ almost surely; each $g^J_{\cv\md}$ has
zero integral by the previous paragraph, so the second part gives
$\lvl{\cv} = \tilde{c}_{\cv}$ and
$\sum_{\lev\tsl}(\btheta-\tilde\btheta)_{\cv\md,\lev\tsl}\wav{\lev}{\tsl} = 0$
almost everywhere, whence $\btheta = \tilde\btheta$ by orthonormality. For the
Gram matrix, $\E\norm{\bZ}_2^2 < \infty$ because $\E\norm{\bX}_2^2 < \infty$
and each wavelet is bounded, and for $a \ne 0$,
$a\T\Gram a = \E\{(a\T\bZ)^2\} = \norm{f_a}^2_{L_2(P)} > 0$ by injectivity,
where $f_a$ is the regression function with parameter $a$.

\emph{Nullity when the scaling column is kept.} With $\phi_{00}\equiv 1$ the
column $X_{\cv}\phi_{00}(U_{\md})$ coincides with the unpenalized column
$X_{\cv}$, for each of the $pq$ pairs, so the $pq$ vectors that oppose one to
the other lie in the kernel and are linearly independent. There are no others:
a kernel vector defines
$\sum_{\cv}X_{\cv}\{a_{\cv} + \sum_{\md}(\alpha_{\cv\md} +
\sum_{\lev\tsl}a_{\cv\md,\lev\tsl}\wav{\lev}{\tsl}(U_{\md}))\} = 0$ almost
surely; by Proposition 1 applied to the part with zero integral,
$a_{\cv\md,\lev\tsl} = 0$ and $a_{\cv}+\sum_{\md}\alpha_{\cv\md} = 0$, which
is the space of dimension $pq$ spanned by those vectors. \qed

\begin{remark}[The interval basis]
\label{rem:cdv}
With the basis of \citet*{cohen1993wavelets} the wavelets of level
$\lev \ge \jzero$ are orthogonal to $V_{\jzero}$, which contains the
constants, so they still have zero integral and the penalized part of each
block satisfies (2.2) with no restriction. The $2^{\jzero}$ scaling functions
do not: the constant is the combination $1 = \sum_{\tsl}\mu_{\tsl}
\phi^{\rm int}_{\jzero\tsl}$ with
$\mu_{\tsl} = \int_0^1\phi^{\rm int}_{\jzero\tsl}$, and the subspace with zero
integral is $\{\mu\T\alpha = 0\}$. Keeping the $2^{\jzero}$ scaling columns of
each block unrestricted repeats the constant and makes $\Gram$ singular with
nullity $pq$, as in the last part of Lemma 1; reparametrizing
$\alpha_{\cv\md} = Q\gamma_{\cv\md}$ with $Q$ an orthonormal basis of
$\{\mu\T\alpha = 0\}$ restores injectivity, at the price of
$pq(2^{\jzero}-1)$ unpenalized parameters.
\end{remark}

%%%%%%%%%%%%%%%%%%%%%%%%%%%%%%%%%%%%%%%%%%%%%%%%%%%%%%%%%%%%%%%%%%%%%%%%%%%%%%
% derivations/02-aproximacao-besov.tex
\section{Approximation by the wavelet sieve}
\setcounter{equation}{0}

\subsection*{Proof of Lemma 2}

Fix a pair $(\cv,\md)$ and write $g = \gcomp{\cv}{\md}$ and
$\theta_{\lev\tsl} = \coefs{\cv}{\md}{\lev}{\tsl}$.

\emph{Tail sum.} By Step 2 below the coefficients are square summable, so
$g \in L_2[0,1]$ and, the system being complete and $g$ orthogonal to the
constant, $g = \sum_{\lev\ge0}\sum_{\tsl}\theta_{\lev\tsl}\wav{\lev}{\tsl}$ in
$L_2$. The orthogonal projection onto $\WJ$ retains the levels below $J$, and
Parseval's identity gives
$\norm{g-\PiJ g}^2_{L_2[0,1]} = \sum_{\lev\ge J}
\norm{\theta_{\lev\cdot}}_2^2$.

\emph{From $\ell_\pi$ to $\ell_2$ inside a level.} In $\R^{2^{\lev}}$ one has
$\norm{a}_2 \le \norm{a}_{\pi}$ when $\pi \le 2$, and
$\norm{a}_2 \le (2^{\lev})^{1/2-1/\pi}\norm{a}_{\pi}$ by Hölder's inequality
when $\pi > 2$. In both cases $\norm{\theta_{\lev\cdot}}_2 \le
2^{\lev(1/2-1/\pi)_{+}}\norm{\theta_{\lev\cdot}}_{\pi}$, and
Assumption 5 gives
\[
  \norm{\theta_{\lev\cdot}}_2 \;\le\;
  C_g\,2^{-\lev\{s+1/2-1/\pi-(1/2-1/\pi)_{+}\}}
  \;=\; C_g\,2^{-\lev\seff},
\]
since the exponent equals $s$ when $\pi \ge 2$ and $s-(1/\pi-1/2)$ when
$\pi < 2$. This is where the loss $(1/\pi-1/2)_{+}$ appears.

\emph{Geometric sum and density.}
$\sum_{\lev\ge J}C_g^2 2^{-2\lev\seff} = C_g^2 2^{-2J\seff}/(1-2^{-2\seff})$,
finite because $\seff > 0$, which is the first display of Lemma 2. For the
second, write $f-f_J = \sum_{\cv,\md}X_{\cv}h_{\cv\md}(U_{\md})$ with
$h_{\cv\md} = \gcomp{\cv}{\md}-\PiJ\gcomp{\cv}{\md}$, the levels cancelling.
By the Cauchy-Schwarz inequality over the $pq$ terms and
$|X_{\cv}| \le B_X$,
\[
  \E\{(f-f_J)^2\} \;\le\; pq\sum_{\cv,\md}
  \E\{X_{\cv}^2h_{\cv\md}^2(U_{\md})\}
  \;\le\; pq\,B_X^2\sum_{\cv,\md}\E\{h_{\cv\md}^2(U_{\md})\},
\]
and $\E\{h^2(U_{\md})\} = \int_0^1 h^2 p_{U_{\md}} \le C_U\norm{h}^2_{L_2[0,1]}$,
the marginal density being bounded by $C_U$ because integrating the joint
density over the other $q-1$ coordinates of $[0,1]^{q-1}$ does not increase
the bound. One factor $pq$ comes from the Cauchy-Schwarz inequality and the
other from the sum. \qed

The uniform version is used in Section S5 to bound the empirical norm of the
bias without an appeal to Markov's inequality.

\begin{lemma}[Uniform approximation error]
\label{lem:sup}
Under Assumptions 4 and 5 with $s > 1/\pi$, each component has a continuous
version and, for every $J \ge 0$,
\[
  \norm{\gcomp{\cv}{\md}-\PiJ\gcomp{\cv}{\md}}_{\infty}
  \;\le\; \frac{(L-1)\norm{\psi}_{\infty}C_g}{1-2^{-(s-1/\pi)}}\,
  2^{-J(s-1/\pi)} ,
\]
and therefore
$\normn{\bb} \le \norm{f-f_J}_{\infty} \le pq B_X(L-1)\norm{\psi}_{\infty}
C_g\,2^{-J(s-1/\pi)}/\{1-2^{-(s-1/\pi)}\}$ for every sample.
\end{lemma}

\begin{proof}
At a point $u$ and a level $\lev$ at most $L-1$ periodized translates are
nonzero, because $\psi(2^{\lev}u-\tsl)\ne0$ requires
$2^{\lev}u-\tsl \in (0,L-1)$, so
$\sum_{\tsl}|\wav{\lev}{\tsl}(u)| \le (L-1)2^{\lev/2}\norm{\psi}_{\infty}$.
With $\max_{\tsl}|\theta_{\lev\tsl}| \le \norm{\theta_{\lev\cdot}}_{\pi} \le
C_g2^{-\lev(s+1/2-1/\pi)}$ the level contributes at most
$(L-1)\norm{\psi}_{\infty}C_g2^{-\lev(s-1/\pi)}$ in the supremum norm. Since
$s > 1/\pi$ the series converges absolutely and uniformly, so its sum is
continuous and agrees almost everywhere with $g$; the tail from level $J$ on
is the geometric sum of the displayed bounds, and the statement for $f-f_J$ is
the triangle inequality over the $pq$ terms.
\end{proof}

The next proposition is the one referred to in Section 2.2 of the paper: it
quantifies what the periodized basis charges for a component that does not
match at the two ends of the interval.

\begin{proposition}[Cost of periodization]
\label{prop:periodization}
Let the basis be as in Assumption 4 and let
$g \in \Besov{s}{\pi}{r}[0,1]$, the Besov space of the interval and not of the
circle, with $s > 1/\pi$, $N > s$ and $\int_0^1 g = 0$. Then there is a
constant $C$ depending on $g$, on $\psi$ and on $(s,\pi)$ such that, for every
$J$ with $2^J \ge L-1$,
\[
  \norm{g-\PiJ g}^2_{L_2[0,1]} \;\le\; C\,2^{-2J\seff}
  \;+\; 2(L-1)\norm{\psi}_1^2\norm{g}^2_{\infty}\,2^{-J},
\]
so that the $L_2$ rate is $2^{-J\min(\seff,1/2)}$. If $g$ is continuous with
$g(0) \ne g(1)$, the second term is of the exact order $2^{-J}$ and
$\norm{g-\PiJ g}_{\infty}$ does not converge to zero.
\end{proposition}

\begin{proof}
Let $\tilde g$ be the one periodic extension of $g$, so that
$\theta_{\lev\tsl} = \int_{\R}\tilde g(y)2^{\lev/2}\psi(2^{\lev}y-\tsl)\,dy$.
At a level with $2^{\lev} \ge L-1$, separate the interior translates, whose
support lies inside $[0,1]$, from the at most $L-1$ boundary translates that
cross the identification of $0$ with $1$. For an interior translate,
$\theta_{\lev\tsl} = \inner{Eg}{\psi_{\lev\tsl}}_{\R}$ for any extension $Eg$
of $g$ to the line; taking $E$ a bounded extension operator from
$\Besov{s}{\pi}{r}[0,1]$ to $\Besov{s}{\pi}{r}(\R)$
\citep[Chapter 3]{Triebel-1983} and using the wavelet characterization of the
latter, which for the implication from the function to the sequence requires
only $N > s$ vanishing moments, gives
$\norm{(\theta_{\lev\tsl})_{\rm int}}_{\pi} \le C_12^{-\lev(s+1/2-1/\pi)}$,
whose tail sum is bounded as in Lemma 2. For a boundary translate,
$s > 1/\pi$ makes $g$ bounded and
$|\theta_{\lev\tsl}| \le \norm{g}_{\infty}\norm{\psi_{\lev\tsl}}_1 =
\norm{g}_{\infty}\norm{\psi}_12^{-\lev/2}$, and there are at most $L-1$ of
them per level, so their tail sum is at most
$2(L-1)\norm{\psi}_1^2\norm{g}_{\infty}^22^{-J}$. For the exact order, if $g$
is continuous with $g(0) \ne g(1)$ then on the support of a boundary translate
$\tilde g$ equals $g(1)+o(1)$ to the left of the identification point and
$g(0)+o(1)$ to the right, so
$\theta_{\lev\tsl} = 2^{-\lev/2}\{g(0)-g(1)\}\int_{y>a_{\tsl}}\psi(y)\,dy +
o(2^{-\lev/2})$; since $\int\psi = 0$ and $\psi$ vanishes on no subinterval of
its support, that partial integral is nonzero for some translate, and the
boundary contribution of the level is of order $2^{-\lev}$. The failure of
uniform convergence is that of any wavelet series at a jump.
\end{proof}

Under Assumption 5 nothing changes, since that assumption is already a
periodic one; what Proposition S3.1 gives is the price of its failure. With
the basis of \citet*{cohen1993wavelets} the space $\VJ$ reproduces polynomials
of degree below $N$ and its basis characterizes $\Besov{s}{\pi}{r}[0,1]$
without periodicity \citep[Section 3.9]{Cohen-2003}, so Lemma 2 holds with
that space and the same proof.

%%%%%%%%%%%%%%%%%%%%%%%%%%%%%%%%%%%%%%%%%%%%%%%%%%%%%%%%%%%%%%%%%%%%%%%%%%%%%%
% derivations/03-desenho-produtos.tex
\section{The Gram matrix of the product design}
\setcounter{equation}{0}

Recall $\bPsi(u) \in \R^{K}$, $K = 1+q\NJ$, the vector that collects the
constant and the wavelets of the $q$ modulators, and the identity
$\bZ = \Pi\{\bX \otimes \bPsi(\bU)\}$, where $\Pi$ is the fixed permutation
that moves the $p$ entries $X_{\cv}\cdot 1$ to the front. Since $\Pi$ is
orthogonal, every statement about eigenvalues is invariant to the column
order. Assumption 4 also gives a constant $C_\psi$, depending only on the
family, with
\begin{equation}
  \label{eq:pointwise}
  \sup_{u\in[0,1]}\sum_{\lev<J}\sum_{\tsl}\wav{\lev}{\tsl}(u)^2
  \;\le\; C_\psi\,2^{J} \qquad \text{for every } J \ge 1 :
\end{equation}
at each level at most $L$ periodized wavelets are nonzero at a given $u$ and
$\norm{\wav{\lev}{\tsl}}_{\infty} = 2^{\lev/2}\norm{\psi}_{\infty}$, so the
left side is at most
$L\norm{\psi}^2_{\infty}\sum_{\lev<J}2^{\lev} \le L\norm{\psi}^2_{\infty}2^J$.

\subsection*{Proof of Proposition 2}

\emph{Part (i), the modulators.} For $a \in \R^{K}$ write
$h_a(u) = a\T\bPsi(u) = a_0 + \sum_{\md}h_{a,\md}(u_{\md})$ with
$h_{a,\md} = \sum_{\lev\tsl}a_{\md,\lev\tsl}\wav{\lev}{\tsl}$. The functions
$1$ and $\wav{\lev}{\tsl}(u_{\md})$, $\md = 1,\dots,q$, are orthonormal in
$L_2([0,1]^q,du)$: within one coordinate this is the orthonormality of the
periodized system, and across coordinates
$\int\wav{\lev}{\tsl}(u_{\md})\wav{\lev'}{\tsl'}(u_{\md'})\,du =
\int_0^1\wav{\lev}{\tsl}\int_0^1\wav{\lev'}{\tsl'} = 0$ for
$\md \ne \md'$, because every wavelet integrates to zero. Hence
$\int_{[0,1]^q}h_a^2 = \norm{a}_2^2$ and, by Assumption 3,
$c_U\norm{a}_2^2 \le \E\{h_a(\bU)^2\} \le C_U\norm{a}_2^2$, which is the claim
since $a\T\Gram_{\bPsi}a = \E\{h_a(\bU)^2\}$.

\emph{Part (i), the full matrix.} Let $\delta \in \R^{pK}$ and reshape it as
$\Delta \in \R^{p\times K}$, row $\cv$ holding the coefficients attached to
$X_{\cv}$, so that $\{\bX\otimes\bPsi(\bU)\}\T\delta = \bX\T\Delta\bPsi(\bU)$.
By the tower property,
\[
  \delta\T\Pi\T\Gram\Pi\delta
  = \E\bigl[\{\bX\T\Delta\bPsi(\bU)\}^2\bigr]
  = \E\bigl[\bPsi(\bU)\T\Delta\T\,\E(\bX\bX\T\mid\bU)\,\Delta\bPsi(\bU)\bigr].
\]
Assumption 2 bounds the inner quadratic form between
$\kappa_1\norm{\Delta\bPsi(\bU)}_2^2$ and $\kappa_2\norm{\Delta\bPsi(\bU)}_2^2$
almost surely, and
$\E\norm{\Delta\bPsi(\bU)}_2^2 = \delta\T(I_p\otimes\Gram_{\bPsi})\delta$ lies
between $\lmin(\Gram_{\bPsi})\norm{\delta}_2^2$ and
$\lmaxx(\Gram_{\bPsi})\norm{\delta}_2^2$. Combining the two and using the
first part gives $\gamma = \kappa_1c_U \le \lmin(\Gram)$ and
$\lmaxx(\Gram) \le \kappa_2C_U$. If $\bX$ and $\bU$ are independent, the
conditional second moment matrix is constant, the same computation gives
$\delta\T\Pi\T\Gram\Pi\delta = \delta\T\{\E(\bX\bX\T)\otimes\Gram_{\bPsi}\}
\delta$ for every $\delta$, and two symmetric matrices with the same quadratic
form are equal; the eigenvalues of a Kronecker product of symmetric matrices
are the products of the eigenvalues, and
$\lmin\{\E(\bX\bX\T)\} \ge \kappa_1$ by Jensen's inequality, since $\lmin$ is
concave and $\E(\bX\bX\T) = \E\{\E(\bX\bX\T\mid\bU)\}$.

\emph{Part (ii), the empirical matrix.} First the range:
$\norm{\bZ_i}_2^2 = \norm{\bX_i}_2^2\norm{\bPsi(\bU_i)}_2^2 \le
pB_X^2(1+qC_\psi2^J) = R_J$ by \eqref{eq:pointwise} applied to each of the $q$
coordinates. Set $A_i = (\bZ_i\bZ_i\T-\Gram)/n$, independent, centered and
symmetric, with $\Gramhat - \Gram = \sum_i A_i$. Then
$\opnorm{\bZ_i\bZ_i\T} = \norm{\bZ_i}_2^2 \le R_J$ and
$\opnorm{\Gram} \le R_J$, so $\opnorm{A_i} \le 2R_J/n$; and
$\E\{(\bZ\bZ\T-\Gram)^2\} = \E(\norm{\bZ}_2^2\bZ\bZ\T) - \Gram^2 \preceq
R_J\Gram$, so that
$\opnorm{\sum_i\E A_i^2} \le R_J\lmaxx(\Gram)/n \le R_J\kappa_2C_U/n$. The
matrix Bernstein inequality \citep[Theorem 1.4]{tropp2012user} applied to
$\pm\sum_iA_i$ gives, for every $t>0$,
\[
  \Prob\bigl(\opnorm{\Gramhat-\Gram} \ge t\bigr)
  \;\le\; 2D\exp\Bigl\{-\frac{nt^2}{2R_J(\kappa_2C_U+2t/3)}\Bigr\},
\]
and $t = \gamma/2$ with Weyl's inequality gives
$\lmin(\Gramhat) \ge \lmin(\Gram)-\gamma/2 \ge \gamma/2$ on the complementary
event, which is the display of part (ii). The exponent is $-cn/R_J$ with $c$
fixed and $R_J \asymp pq2^J$, while $\log(2D) \asymp \log(pq2^J)$, so the
bound tends to zero exactly when $n/(pq2^J) - \log(pq2^J) \to \infty$, which
is implied by $pq2^J\log(pq2^J)/n \to 0$. On that event
$\compat^2(S;\Gramhat) \ge \gamma/2$ for every support, by
\eqref{eq:compat-lmin}. \qed

\begin{remark}[What the two assumptions exclude]
\label{rem:excl}
Assumption 2 fails exactly when a nontrivial linear combination of the linear
covariates is almost surely a function of the modulators on a set of positive
probability; in the canonical case $X_1\equiv1$ and $X_2 = h(U_1)$ the
conditional second moment matrix has rank one and the population Gram matrix
degenerates as $J$ grows, because $h(u_1)\wav{\lev}{\tsl}(u_1)$ is approached
by the span of the basis. In Assumption 3 it is the joint density that is
bounded below: if $U_2 = U_1$ almost surely, both marginals are uniform but
$\wav{\lev}{\tsl}(U_1)-\wav{\lev}{\tsl}(U_2) \equiv 0$ and
$\Gram_{\bPsi}$ is singular. Independence between $\bX$ and $\bU$ buys the
exact factorization of part (i) and nothing else.
\end{remark}

%%%%%%%%%%%%%%%%%%%%%%%%%%%%%%%%%%%%%%%%%%%%%%%%%%%%%%%%%%%%%%%%%%%%%%%%%%%%%%
% derivations/04-oraculo.tex
\section{The oracle inequality}
\setcounter{equation}{0}

Four lemmas precede the proof of Theorem 2.

\begin{lemma}[Profiling and the Schur complement]
\label{lem:profile}
Suppose $\Amat$ has rank $p$. Then $\thetahat$ minimizes
$\normn{\MA\bY-\Btil\btheta}^2+2\pen\norm{\btheta}_1$ and
$\chat = (\Amat\T\Amat)^{-1}\Amat\T(\bY-\Bmat\thetahat)$; with
$v = \thetahat-\thetastar$,
\begin{equation}
  \label{eq:decomp}
  \widehat{f}-f_J \;=\; \Btil v + \PA(\bb+\bveps);
\end{equation}
$\GBA$ is the Schur complement of the block $\Amat$ in $\Gramhat$ and
\begin{equation}
  \label{eq:schur}
  \gtil \;=\; \lmin(\GBA) \;\ge\; \lmin(\Gramhat);
\end{equation}
and if $\lmin(\Gramhat)>0$ the minimizer in (2.5) is unique.
\end{lemma}

\begin{proof}
The objective is convex and, for fixed $\btheta$, quadratic and coercive in
$\bc$, with minimum at the stated value and minimized value
$\normn{\MA(\bY-\Bmat\btheta)}^2 = \normn{\MA\bY-\Btil\btheta}^2$, since $\MA$
is a symmetric idempotent; minimizing first in $\bc$ and then in $\btheta$
gives the same set of minimizers as minimizing jointly. For
\eqref{eq:decomp}, $\Amat\chat = \PA(\bY-\Bmat\thetahat)$ and
$\Amat\bc = \PA\Amat\bc$, and $\bY = \Amat\bc+\Bmat\thetastar+\bb+\bveps$, so
$\Amat(\chat-\bc) = \PA(\bb+\bveps)-\PA\Bmat v$ and
$\widehat{f}-f_J = \MA\Bmat v + \PA(\bb+\bveps)$. For \eqref{eq:schur}, the
variational characterization of the Schur complement gives, for every
$v \in \R^d$,
\[
  v\T\GBA v \;=\; \min_{u\in\R^p}(u\T,v\T)\,\Gramhat\,(u\T,v\T)\T
  \;\ge\; \min_{u}\lmin(\Gramhat)(\norm{u}_2^2+\norm{v}_2^2)
  \;=\; \lmin(\Gramhat)\norm{v}_2^2 .
\]
Finally $\lmin(\Gramhat)>0$ makes the quadratic term strictly convex, and
adding a convex function preserves strict convexity.
\end{proof}

\begin{lemma}[Calibration of the penalty]
\label{lem:calib}
Under Assumption 1, for every $\alpha \in (0,1)$,
$\Prob(\norm{\Btil\T\bveps/n}_{\infty} > \pen_0) \le \alpha$ with
$\pen_0 = \sigma\smax\{2\log(2d/\alpha)/n\}^{1/2}$.
\end{lemma}

\begin{proof}
Condition on the covariates, so that $\Btil$ and $\smax$ are fixed and the
errors are independent and sub-Gaussian with parameter $\sigma$. For a column
$a$, $(\Btil\T\bveps)_a/n$ is sub-Gaussian with parameter
$\sigma\lVert\Btil^{(a)}\rVert_2/n$, and, $\MA$ being an orthogonal
projection, $\lVert\Btil^{(a)}\rVert_2 \le \lVert\Bmat^{(a)}\rVert_2 =
\sqrt{n}\,(\Gramhat)_{aa}^{1/2} \le \sqrt{n}\,\smax$. With
$t = \pen_0$ the sub-Gaussian tail bound is at most $\alpha/d$ per column; a
union bound over the $d$ columns and the tower property finish the proof.
\end{proof}

\begin{lemma}[Empirical column norms]
\label{lem:columns}
Under Assumptions 2 to 4, every penalized column has
$\E\{(\Gramhat)_{aa}\} \le B_X^2C_U$ and, for every $t>0$, with
$R'_J = B_X^2\norm{\psi}^2_{\infty}2^{J-1}$,
\[
  \Prob(\smax^2 > B_X^2C_U+t) \;\le\;
  d\exp\Bigl\{-\frac{nt^2}{2(R'_JB_X^2C_U+R'_Jt/3)}\Bigr\} .
\]
In particular $\Prob(\smax^2 \le 2B_X^2C_U) \to 1$ whenever
$2^J\log d/n \to 0$.
\end{lemma}

\begin{proof}
For the column $a = (\cv\md,\lev\tsl)$ the summand is
$W_i = X_{i\cv}^2\wav{\lev}{\tsl}(U_{i\md})^2 \ge 0$, with
$\E(W) \le B_X^2\int_0^1\wav{\lev}{\tsl}^2p_{U_{\md}} \le B_X^2C_U$ by
orthonormality and the bounded marginal density, and
$0 \le W_i \le B_X^2\norm{\psi}_{\infty}^22^{J-1} = R'_J$ because
$\lev \le J-1$. Then $\Var(W) \le \E(W^2) \le R'_J\E(W) \le R'_JB_X^2C_U$, and
the one sided Bernstein inequality with a union bound over the $d$ columns
gives the display. With $t = B_X^2C_U$ the exponent is
$-3nB_X^2C_U/(8R'_J)$ and $R'_J \asymp 2^J$.
\end{proof}

\begin{lemma}[The $\ell_1$ norm of the sieve oracle]
\label{lem:l1}
Under Assumption 5, for every pair $(\cv,\md)$ and every level,
$\sum_{\tsl}|\coefs{\cv}{\md}{\lev}{\tsl}| \le
2^{\lev(1-1/\pi)}\norm{\theta^{*}_{\cv\md,\lev\cdot}}_{\pi} \le
C_g2^{\lev(1/2-s)}$, with no dependence on $\pi$, and therefore
$\norm{\thetastar}_1 \le pqC_g\sum_{\lev<J}2^{\lev(1/2-s)}$, which is
$O(1)$ if $s>1/2$, of order $J$ if $s = 1/2$ and of order
$2^{J(1/2-s)}$ if $s<1/2$.
\end{lemma}

\begin{proof}
Hölder's inequality in $\R^{2^{\lev}}$ gives
$\norm{a}_1 \le (2^{\lev})^{1-1/\pi}\norm{a}_{\pi}$, and Assumption 5 bounds
the $\ell_\pi$ norm of the level by $C_g2^{-\lev(s+1/2-1/\pi)}$; the resulting
exponent is $\lev(1-1/\pi)-\lev(s+1/2-1/\pi) = \lev(1/2-s)$, in which $\pi$
has disappeared. Summing over the levels below $J$ and over the $pq$ blocks
gives the three regimes of the geometric sum.
\end{proof}

\subsection*{Proof of Theorem 2}

\emph{Step 1, the basic inequality.} By Lemma S5.1, $\thetahat$ minimizes the
residualized objective; comparing its value at $\thetahat$ with its value at
$\thetastar$, and using $\MA\Amat = 0$, so that
$\MA\bY = \Btil\thetastar + \MA(\bb+\bveps)$, gives, with
$r = \MA(\bb+\bveps)$,
\begin{equation}
  \label{eq:basic}
  \normn{\Btil v}^2 + 2\pen\norm{\thetahat}_1
  \;\le\; 2\inner{r}{\Btil v}_n + 2\pen\norm{\thetastar}_1 .
\end{equation}

\emph{Step 2, the two noises.} Since $\MA\Btil = \Btil$ and $\MA$ is
symmetric, $\inner{r}{\Btil v}_n = \inner{\bb+\bveps}{\Btil v}_n$. The
stochastic part is handled by duality: on $\mathcal{T}$,
$2|\inner{\bveps}{\Btil v}_n| \le
2\norm{\Btil\T\bveps/n}_{\infty}\norm{v}_1 \le \pen\norm{v}_1$. The
deterministic part is handled by the Cauchy-Schwarz and Young inequalities,
and not by duality, because $\bb$ is not small coordinatewise:
$2|\inner{\bb}{\Btil v}_n| \le 2\normn{\bb}\normn{\Btil v} \le
\normn{\Btil v}^2/2 + 2\normn{\bb}^2$. Substituting both in
\eqref{eq:basic},
\begin{equation}
  \label{eq:basic2}
  \tfrac12\normn{\Btil v}^2 + 2\pen\norm{\thetahat}_1
  \;\le\; \pen\norm{v}_1 + 2\pen\norm{\thetastar}_1 + 2\normn{\bb}^2 .
\end{equation}

\emph{Step 3, the slow rate.} Bounding
$\norm{v}_1 \le \norm{\thetahat}_1+\norm{\thetastar}_1$ in
\eqref{eq:basic2} gives
$\tfrac12\normn{\Btil v}^2 + \pen\norm{\thetahat}_1 \le
3\pen\norm{\thetastar}_1 + 2\normn{\bb}^2$, which is part (i). No property of
$\Gramhat$ was used.

\emph{Step 4, the cone.} Since $\thetastar$ vanishes off $\Szero$,
$\norm{\thetahat}_1 \ge \norm{\thetastar}_1 - \norm{v_{\Szero}}_1 +
\norm{v_{\Szero^c}}_1$ and
$\norm{v}_1 = \norm{v_{\Szero}}_1+\norm{v_{\Szero^c}}_1$; substituting in
\eqref{eq:basic2} and cancelling $2\pen\norm{\thetastar}_1$,
\begin{equation}
  \label{eq:cone}
  \tfrac12\normn{\Btil v}^2 + \pen\norm{v_{\Szero^c}}_1
  \;\le\; 3\pen\norm{v_{\Szero}}_1 + 2\normn{\bb}^2 .
\end{equation}

\emph{Step 5, two cases.} If $2\normn{\bb}^2 \le \pen\norm{v_{\Szero}}_1$,
then \eqref{eq:cone} gives
$\tfrac12\normn{\Btil v}^2 + \pen\norm{v_{\Szero^c}}_1 \le
4\pen\norm{v_{\Szero}}_1$; by the Cauchy-Schwarz inequality and
\eqref{eq:schur}, $\norm{v_{\Szero}}_1 \le \sqrt{s_0}\norm{v}_2 \le
\sqrt{s_0}\normn{\Btil v}/\sqrt{\gtil}$, so
$\normn{\Btil v} \le 8\pen\sqrt{s_0/\gtil}$ and
$\norm{v}_2^2 \le \normn{\Btil v}^2/\gtil \le 64\pen^2s_0/\gtil^2$; the same
display gives $\norm{v_{\Szero^c}}_1 \le 4\norm{v_{\Szero}}_1$ and hence
$\norm{v}_1 \le 5\norm{v_{\Szero}}_1 \le 40\pen s_0/\gtil$. If instead
$2\normn{\bb}^2 > \pen\norm{v_{\Szero}}_1$, then \eqref{eq:cone} gives
$\tfrac12\normn{\Btil v}^2 < 8\normn{\bb}^2$, that is
$\normn{\Btil v}^2 < 16\normn{\bb}^2$, together with
$\pen\norm{v}_1 < 10\normn{\bb}^2$ and
$\norm{v}_2^2 < 16\normn{\bb}^2/\gtil$. Each quantity is bounded by the
maximum of the two cases, hence by their sum, which is part (ii).

\emph{Step 6, prediction.} By \eqref{eq:decomp},
$\widehat{f}-f = \Btil v + \PA(\bb+\bveps) - \bb$, and $\PA$ being a
projection, $\normn{\PA\bb} \le \normn{\bb}$, which gives the triangle
inequality of part (iii). Conditionally on the covariates the errors are
independent, centered and of variance at most $\sigma^2$, so
\[
  \E(\normn{\PA\bveps}^2\mid\bX,\bU)
  = \frac1n\tr\{\PA\,\mathrm{Cov}(\bveps\mid\bX,\bU)\}
  \le \frac{\sigma^2}{n}\tr(\PA) \le \frac{\sigma^2p}{n},
\]
since the diagonal entries of $\PA$ are nonnegative and its trace is the rank
of $\Amat$. \qed

\begin{remark}[Why there is no cone condition]
\label{rem:nocone}
The usual treatment of unpenalized coordinates \citep[Section
6.9]{buhlmann2011statistics} places them inside the support and requires
compatibility on the cone $\norm{\delta_{S^c}}_1 \le 3\norm{\delta_S}_1$.
Profiling removes them from the problem, and \eqref{eq:schur} transfers the
full smallest eigenvalue of $\Gramhat$ to the residualized block, so that
Proposition 2(ii) is used as it stands and the constant does not depend on
$\Szero$. If only the restricted version of the design condition were
available, the cone would return, and with the constant $4$ rather than $3$,
since Step 5 produces $\norm{v_{\Szero^c}}_1 \le 4\norm{v_{\Szero}}_1$: the
bias consumes part of the slack.
\end{remark}

\begin{corollary}[The two rates with the design condition instantiated]
\label{cor:instant}
Add Assumptions 2 to 5 and suppose $2^J\log d/n \to 0$, which is implied by
the requirement of Proposition 2(ii). Take
$\pen = 2\sigma B_X\sqrt{2C_U}\{2\log(2d/\alpha)/n\}^{1/2}$. Then, with
probability at least $1-\alpha-\alpha'-\alpha''$ up to terms that tend to
zero,
\[
  \normn{\widehat{f}-f}^2 \;\lesssim\;
  \frac{\sigma^2B_X^2C_U}{\kappa_1c_U}\,\frac{s_0\log(pq\NJ)}{n}
  \;+\; \frac{\mathcal{B}^2_J}{\alpha'} \;+\; \frac{\sigma^2p}{n\alpha''},
\]
and, with no condition on the design, the same bound with
$\pen\norm{\thetastar}_1$ of Lemma S5.4 in place of the first term.
\end{corollary}

\begin{proof}
Intersect four events: $\mathcal{T}$, of probability at least $1-\alpha$, by
Lemma S5.2; $\{\smax^2 \le 2B_X^2C_U\}$, by Lemma S5.3, which makes the stated
$\pen$ admissible, that is $\pen \ge 2\pen_0$;
$\{\lmin(\Gramhat)\ge\gamma/2\}$, by Proposition 2(ii), which with
\eqref{eq:schur} gives $\gtil \ge \gamma/2$; and
$\{\normn{\bb}^2 \le \mathcal{B}_J^2/\alpha'\}$, which is Markov's inequality
applied to $\E\normn{\bb}^2 = \norm{f-f_J}^2_{L_2(P)} \le \mathcal{B}_J^2$,
the identity holding because the observations are identically distributed.
Substituting in parts (ii) and (iii) of Theorem 2 gives the statement, and
part (i) with Lemma S5.4 gives the unconditional version.
\end{proof}

%%%%%%%%%%%%%%%%%%%%%%%%%%%%%%%%%%%%%%%%%%%%%%%%%%%%%%%%%%%%%%%%%%%%%%%%%%%%%%
% derivations/05-taxas.tex
\section{Rates}
\setcounter{equation}{0}

Throughout this section $p$, $q$, $\sigma$, $B_X$, $c_U$, $C_U$, $\kappa_1$,
$\kappa_2$, $C_g$, $s$, $\pi$, $r$ and $\alpha$ are fixed as $n$ grows,
$J = J_n \to \infty$, $d_n = pq(2^{J_n}-1)$,
$L_n = \log(2d_n/\alpha) \asymp J_n$, and $\penn$ is as in (3.1). Since the
constants of Theorem 2 depend on the confidence levels only through
$1/\alpha'$ and $1/\alpha''$, taking $\alpha' = \alpha'' = \epsilon/3$ gives,
for each $\epsilon>0$, a constant $M_\epsilon$ with
$\Prob(\normn{\widehat{f}-f}^2 > M_\epsilon r_n) < \epsilon$, which is the
meaning of the statements in $O_p$ below. We also write
$\Lambda = \kappa_2C_U+\gamma/2$.

The first lemma is Theorem 2 read at an arbitrary comparator. It consumes no
new step: the proof of that theorem uses the oracle only through the
decomposition of $\bY$ and through the size of its support.

\begin{lemma}[Theorem 2 with an arbitrary comparator]
\label{lem:comparator}
Under Assumption 1, with $\Amat$ of rank $p$ and $\pen \ge 2\pen_0$, on the
event $\mathcal{T}$ the following hold simultaneously for every
$\bar\bbeta = (\bar\bc,\thbar)$, with $\bbar = f - \bZ\bar\bbeta$,
$\bar M = |\supp(\thbar)|$ and $\bar v = \thetahat-\thbar$:
$\tfrac12\normn{\Btil\bar v}^2 + \pen\norm{\thetahat}_1 \le
3\pen\norm{\thbar}_1 + 2\normn{\bbar}^2$; if $\gtil>0$,
$\normn{\Btil\bar v}^2 \le 64\pen^2\bar M/\gtil + 16\normn{\bbar}^2$ and
$\norm{\bar v}_2^2 \le 64\pen^2\bar M/\gtil^2 +
16\normn{\bbar}^2/\gtil$; and
\begin{equation}
  \label{eq:oracle-min}
  \normn{\widehat{f}-f}^2 \;\le\;
  3\min_{\bar\bbeta}\Bigl\{\frac{64\pen^2\bar M}{\gtil}
  + 20\normn{\bbar}^2\Bigr\} + 3\normn{\PA\bveps}^2 .
\end{equation}
\end{lemma}

\begin{proof}
Follow the proof of Theorem 2 with
$(\bar\bc,\thbar,\bbar,\supp(\thbar),\bar M,\bar v)$ in place of
$(\bc,\thetastar,\bb,\Szero,s_0,v)$. Step 1 uses only
$\bY = \Amat\bar\bc+\Bmat\thbar+\bbar+\bveps$, which is the definition of
$\bbar$, and $\MA\Amat = 0$. Step 2 uses $\mathcal{T}$, which does not depend
on the comparator, and $\normn{\MA\bbar} \le \normn{\bbar}$. Steps 3 to 5 are
algebra with $\supp(\thbar)$ in place of $\Szero$, using
$\thbar_{\supp(\thbar)^c} = 0$, the Cauchy-Schwarz inequality and
\eqref{eq:schur}. Step 6 gives
$\widehat{f}-\bZ\bar\bbeta = \Btil\bar v + \PA(\bbar+\bveps)$ and hence
$\widehat{f}-f = \Btil\bar v+\PA(\bbar+\bveps)-\bbar$. Since $\mathcal{T}$,
$\gtil$ and the condition on $\pen$ do not involve $\bar\bbeta$, the
inequalities hold simultaneously on the same event and the minimum may be
taken; \eqref{eq:oracle-min} follows from
$(x+y+z)^2 \le 3(x^2+y^2+z^2)$.
\end{proof}

\subsection*{Proof of Lemma 3}

\emph{The envelope.} Fix $\epsilon>0$ and count
$N(\epsilon) = \#\{a : |\theta^{*}_a| > \epsilon\}$. Inside a block and a
level, Markov's inequality in the $\ell_\pi$ norm gives
$\#\{\tsl : |\coefs{\cv}{\md}{\lev}{\tsl}|>\epsilon\} \le
\epsilon^{-\pi}\norm{\theta^{*}_{\cv\md,\lev\cdot}}^{\pi}_{\pi}$, and
Assumption 5 gives
$\norm{\theta^{*}_{\cv\md,\lev\cdot}}^{\pi}_{\pi} \le
C_g^{\pi}2^{-\lev\pi(s+1/2-1/\pi)} = C_g^{\pi}2^{\lev(1-a)}$ with
$a = \pi(s+1/2)>1$. Since a level has $2^{\lev}$ coordinates,
\[
  N_{\cv\md}(\epsilon) \;\le\; \sum_{\lev\ge0}
  \min\{2^{\lev},\, C_g^{\pi}\epsilon^{-\pi}2^{\lev(1-a)}\} .
\]
The two arguments agree at $2^{\lev_\epsilon a} = (C_g/\epsilon)^{\pi}$, that
is at $2^{\lev_\epsilon} = (C_g/\epsilon)^{\tau}$, since
$\pi/a = (s+1/2)^{-1} = \tau$. Below that level the minimum is $2^{\lev}$ and
the sum is at most $2^{\lev_\epsilon}$; above it the minimum is the second
term, which decreases geometrically with ratio $2^{1-a}<1$, and the sum is
$2^{\lev_\epsilon}/(1-2^{1-a})$. Adding the two parts and the $pq$ blocks
gives $N(\epsilon) \le pqA_{s,\pi}(C_g/\epsilon)^{\tau}$; truncating at
$\lev<J$ only decreases the count, so the bound holds for every $J$. Finally,
$\epsilon < \theta^{*}_{(k)}$ implies $N(\epsilon) \ge k$, hence
$\epsilon \le C_g(pqA_{s,\pi}/k)^{1/\tau}$, and letting $\epsilon$ increase to
$\theta^{*}_{(k)}$ gives the envelope; and $\tau = (s+1/2)^{-1} < 2$ for every
$s>0$.

\emph{The tail.} With $2/\tau>1$,
$\sum_{k>M}\theta^2_{(k)} \le C_\tau^2\sum_{k>M}k^{-2/\tau} \le
C_\tau^2\int_M^{\infty}x^{-2/\tau}dx = C_\tau^2M^{1-2/\tau}\tau/(2-\tau)$.
\qed

We also record that, on the event of Proposition 2(ii) with $t = \gamma/2$,
\begin{equation}
  \label{eq:lmax}
  \normn{\Bmat\delta}^2 \le \lmaxx(\Gramhat)\norm{\delta}_2^2
  \quad\text{and}\quad \lmaxx(\Gramhat) \le \Lambda,
\end{equation}
the first because $\Bmat\T\Bmat/n$ is a principal submatrix of $\Gramhat$ and
the second by Weyl's inequality and Proposition 2(i).

\subsection*{Proof of Proposition 3}

By parts (i) and (iii) of Theorem 2, on $\mathcal{T}$,
$\normn{\Btil v}^2 \le 6\pen\norm{\thetastar}_1+4\normn{\bb}^2$ and
$\normn{\widehat{f}-f}^2 \le 3(\normn{\Btil v}^2+4\normn{\bb}^2+
\normn{\PA\bveps}^2)$, and no property of $\Gramhat$ enters. Lemma S5.3 makes
$\penn$ admissible with probability tending to one, since
$2^{J_n}\log d_n/n \to 0$ in every regime considered; Markov's inequality
applied to $\E\normn{\bb}^2 \le \mathcal{B}^2_{J_n}$, from Lemma 2, and to
$\E(\normn{\PA\bveps}^2\mid\bX,\bU) \le \sigma^2p/n$ gives the general
statement. If $s>1/2$, Lemma S5.4 gives $\norm{\thetastar}_1 = O(1)$ and
$\penn\norm{\thetastar}_1 \asymp (J_n/n)^{1/2}$, and it suffices that the bias
not dominate, that is $2^{-2J_n\seff} \le (J_n/n)^{1/2}$, which holds for
$J_n \ge \log_2(n/J_n)/(4\seff)$. If $s<1/2$, then
$\norm{\thetastar}_1 \asymp 2^{J_n(1/2-s)}$, the two terms are
$(J_n/n)^{1/2}2^{J_n(1/2-s)}$ and $2^{-2J_n\seff}$, and equating them gives
$2^{J_n} \asymp (n/J_n)^{1/(1-2s+4\seff)}$ with common value
$(J_n/n)^{2\seff/(1-2s+4\seff)}$; when $\seff = s$ the denominator is
$1+2s$. \qed

\subsection*{Proof of Theorem 3}

By the definition of $J_n$,
$(n/\log n)^{1/(2\seff+1)} \le 2^{J_n} < 2(n/\log n)^{1/(2\seff+1)}$ and
$J_n \asymp \log n$, so
\[
  \frac{pq2^{J_n}\log(pq2^{J_n})}{n}
  \;\asymp\; \frac{(n/\log n)^{1/(2\seff+1)}\log n}{n}
  \;=\; \Bigl(\frac{\log n}{n}\Bigr)^{2\seff/(2\seff+1)}
  \;\longrightarrow\; 0,
\]
which is the requirement of Proposition 2(ii); \eqref{eq:schur} then gives
$\gtil \ge \gamma/2$ with probability tending to one. On that event
Corollary S5.1 with $s_0 \le d_n$ gives
\[
  \normn{\widehat{f}-f}^2 \;\lesssim\;
  \frac{\sigma^2B_X^2C_U}{\kappa_1c_U}\frac{d_nL_n}{n}
  + \frac{\mathcal{B}^2_{J_n}}{\alpha'} + \frac{\sigma^2p}{n\alpha''} .
\]
With $d_n \asymp pq2^{J_n}$ and $L_n \asymp \log n$ the first term is the
display above; the second is
$\mathcal{B}^2_{J_n} \asymp 2^{-2J_n\seff} \le
(\log n/n)^{2\seff/(2\seff+1)}$ by the lower bound on $2^{J_n}$; and the third
is $O(1/n)$, of strictly smaller order because $2\seff/(2\seff+1)<1$. \qed

\subsection*{Proof of Corollary 1}

Inside a block the wavelets of level below $J$ are orthonormal in $L_2[0,1]$,
so
$\lVert\widehat{g}_{\cv\md}-\PiJ\gcomp{\cv}{\md}\rVert^2_{L_2[0,1]} =
\sum_{\lev\tsl}v^2_{\cv\md,\lev\tsl}$, and summing over the $pq$ blocks gives
$\norm{v}_2^2$; adding the bias of Lemma 2 in each block,
\[
  \sum_{\cv,\md}\lVert\widehat{g}_{\cv\md}-\gcomp{\cv}{\md}
  \rVert^2_{L_2[0,1]} \;\le\; 2\norm{v}_2^2
  + \frac{2pqC_g^2}{1-2^{-2\seff}}2^{-2J\seff} .
\]
Theorem 2(ii) bounds $\norm{v}_2^2$ by
$64\pen^2s_0/\gtil^2+16\normn{\bb}^2/\gtil$, and on the event of Theorem 3 the
quantity $\gtil$ is bounded below by the constant $\gamma/2$, so the two terms
are of the same order as those of Theorem 3, the only difference being a
constant with $\gamma^{-2}$ in place of $\gamma^{-1}$. This is where the
design condition is used a second time. For the levels, Lemma S5.1 gives
$\Amat(\chat-\bc) = \PA(\bb+\bveps)-\PA\Bmat v$, so that
$\normn{\Amat(\chat-\bc)} \le \normn{\bb}+\normn{\PA\bveps}+
\lmaxx(\Gramhat)^{1/2}\norm{v}_2$, while $\Amat\T\Amat/n$ is a principal
submatrix of $\Gramhat$, whence
$\norm{\chat-\bc}_2^2 \le \normn{\Amat(\chat-\bc)}^2/(\gamma/2)$. Finally the
components of a functional coefficient depend on distinct coordinates of $u$
and the basis is orthonormal, so
$\lVert\widehat{\fcoef{\cv}}-\fcoef{\cv}\rVert_{L_2([0,1]^q)} \le
|\widehat{c}_{\cv}-\lvl{\cv}| + \sum_{\md}\lVert\widehat{g}_{\cv\md}-
\gcomp{\cv}{\md}\rVert_{L_2[0,1]}$, with $q$ fixed. The same argument applies
verbatim under the conditions of Theorem 1. \qed

\subsection*{Proof of Theorem 1}

\emph{A nonasymptotic bound.} Apply Lemma S6.1 with
$\bar\bbeta = (\bc,\thstarM{M})$, the best $M$ term approximation of the
oracle. Then $\bar M \le M$ and
$\bbar = \bb + \Bmat(\thetastar-\thstarM{M})$, so that, by the
Cauchy-Schwarz inequality, Lemma 3 and \eqref{eq:lmax},
\[
  \normn{\bbar}^2 \;\le\; 2\normn{\bb}^2
  + 2\lmaxx(\Gramhat)\norm{\thetastar-\thstarM{M}}_2^2
  \;\le\; 2\normn{\bb}^2 + \frac{2\Lambda C_\tau^2\tau}{2-\tau}M^{1-2/\tau} .
\]
Substituting in \eqref{eq:oracle-min}, on the intersection of $\mathcal{T}$
with $\{\lmin(\Gramhat)\ge\gamma/2\}$,
\[
  \normn{\widehat{f}-f}^2 \;\le\; 3\min_{1\le M\le d}h(M)
  + 120\normn{\bb}^2 + 3\normn{\PA\bveps}^2,
  \qquad
  h(M) = \frac{64\pen^2}{\gtil}M
  + \frac{40\Lambda C_\tau^2\tau}{2-\tau}M^{1-2/\tau} .
\]
The function $h$ is strictly convex on the positive half line, since
$1-2/\tau<0$, and $h'$ vanishes at
$\Mstar = \{5\Lambda C_\tau^2\gtil/(8\pen^2)\}^{\tau/2}$, where
$h(\Mstar) = 2(2-\tau)^{-1}(64\pen^2/\gtil)\Mstar$; moreover
$h(\lceil\Mstar\rceil) \le h(\Mstar)+64\pen^2/\gtil \le 2h(\Mstar)$ whenever
$\Mstar \ge 1$, because $(2-\tau)/2<1$.

\emph{The resolution.} With $\penn^2 \asymp \sigma^2B_X^2C_UJ_n/n$,
$\gtil \in [\gamma/2,\Lambda]$ and the other constants fixed,
$\Mstar_n \asymp \penn^{-\tau} \asymp (n/J_n)^{\tau/2}$ and
$h(\Mstar_n) \asymp \penn^2\Mstar_n \asymp (J_n/n)^{(2-\tau)/2}$. Writing
$J_n = \lceil c\log_2 n\rceil$, so that $2^{J_n} \asymp n^{c}$, the three
inequalities of the window (3.2) are, in order: $\Mstar_n \le d_n$, which asks
$(n/J_n)^{\tau/2} \lesssim pq\,n^{c}$, that is $c \ge \tau/2$ up to
logarithms; that the bias not dominate, $n^{-2c\seff} \lesssim
(J_n/n)^{(2-\tau)/2}$, that is $c \ge (2-\tau)/(4\seff)$; and the requirement
of Proposition 2(ii), $n^{c-1}\log n \to 0$, that is $c<1$. The term
$\sigma^2p/n$ is of smaller order because $(2-\tau)/2<1$. The window is
nonempty if and only if $\tau/2<1$, which always holds, and
$(2-\tau)/(4\seff)<1$, that is $\seff>(2-\tau)/4$.

\emph{Under Assumption 5 the hypothesis is automatic.} By Lemma 3 the
envelope holds with $\tau = (s+1/2)^{-1}$, whence $2-\tau = 4s/(2s+1)$,
$(2-\tau)/2 = 2s/(2s+1)$ and $(2-\tau)/4 = s/(2s+1)$, which is the condition
$\seff>s/(2s+1)$ of the statement; for $\pi \ge 2$ it holds for every $s>0$,
and for $\pi<2$ it reads $2s^2/(2s+1)>1/\pi-1/2$. The rate for the components
is the argument of Corollary 1 with
$\norm{\thetahat-\thetastar}_2 \le \norm{\bar v}_2 +
\norm{\thetastar-\thstarM{M}}_2$ and Lemma S6.1, the only difference being an
extra factor $\gtil^{-1}$ in the first term, which does not change the order.
Finally $2s/(2s+1) \ge 2\seff/(2\seff+1)$, with equality only when
$\seff = s$, by monotonicity of $x \mapsto 2x/(2x+1)$ and $\seff \le s$, so
the rate is never worse than that of Theorem 3 and is strictly better when
$\pi<2$. \qed

\begin{remark}[What compressibility buys]
\label{rem:s0}
Read at the oracle, Theorem 2(ii) has $s_0 = d$ and delivers the rate of the
sieve; read at a sparse comparator, through Lemma S6.1, it delivers the rate
of Theorem 1. It is the same inequality in both cases, and neither says
anything about $\supp(\thetahat)$: no minimal signal and no
irrepresentability condition is used, and what is gained is a rate governed by
the effective dimension $\Mstar_n$ rather than by $d_n$.
\end{remark}

%%%%%%%%%%%%%%%%%%%%%%%%%%%%%%%%%%%%%%%%%%%%%%%%%%%%%%%%%%%%%%%%%%%%%%%%%%%%%%
\bibhang=1.7pc
\bibsep=2pt
\fontsize{9}{14pt plus.8pt minus .6pt}\selectfont
\renewcommand\bibname{\large \bf References}
\bibliographystyle{chicago}
\bibliography{references_1}

\end{document}
